%---------------------------------------------------------------------------------
%                西南交通大学研究生学位论文：第三章内容
%---------------------------------------------------------------------------------
\chapter{基于逻辑增强与异构协同的政务数据查询生成方法}

\section{引言}

% TODO: 形式化定义输入为自然语言问题 q、数据库模式 S、辅助知识库 K 与样例库 F，
% 输出为可正确执行的 SQL y。
% 说明本章目标不是单纯生成 SQL，而是在复杂政务语义下提升查询正确性、执行成功率与鲁棒性。
% 给出一个政务问句到 SQL 的简短示例。

% TODO: 介绍政务复杂业务逻辑特征（聚焦方法动机）：
% 中文简称与别名多、指标口径隐含、时间条件省略、多表关联复杂、字段命名与业务术语不一致等。
% 说明这些特征会分别影响模式链接、候选生成与结果判别。

% TODO: 讲述方法设计目标（三点）：
% 1. 提升模式链接准确性
% 2. 提升候选生成多样性与质量
% 3. 提升最终查询结果稳定性


\section{总体框架设计}

% 建议图表：图3-1 逻辑增强与异构协同的 Text-to-SQL 框架图（核心）

% TODO: 用图3-1 把离线与在线两条链路讲清楚。
% 离线阶段：构建 VecDB、领域知识库和动态样例库。
% 在线阶段：逻辑增强模式链接 -> 双路候选生成 -> 查询修复 -> 候选判别。
% 总结各模块之间的数据流转关系，形成"知识准备 - 生成 - 修复 - 选择"的完整闭环。


\section{离线知识准备}

\subsection{向量库构建}

% TODO: 说明向量库存储的对象（表名、字段名、字段注释、样例值、业务描述文本等），
% 采用何种嵌入模型和索引方式。
% 强调它在在线阶段承担相似实体召回、别名匹配和相关样例检索的作用。


\subsection{政务领域知识库构建}

% TODO: 写清知识来源（指标词典、业务术语表、字段解释文档、统计规则和历史问答知识）。
% 说明知识如何归一化组织，例如"术语 - 标准表达 - 对应字段/取值 - 业务说明"的结构。


\subsection{动态样例库构建}

% TODO: 重点写"动态"，即样例不是固定 few-shot，
% 而是根据当前问题检索最相似的历史 NL-SQL 样本。
% 说明样例来源、筛选规则、检索策略与更新机制。


\section{基于逻辑增强的模式链接}

% TODO: 本章第一创新点，重点写。分三层展开：
% 1. 术语归一与别名扩展
% 2. 列值提示与候选字段约束
% 3. LA-Schema 的组织方式
% 给出一个中文政务问题示例，展示原始 Schema 与增强后 Schema 在提示效果上的差异。


\section{异构候选生成}

\subsection{逻辑映射生成器}

% TODO: 讲述如何应用 LA-Schema、检索样例、用户问题以及输出约束。
% 强调优势在于效率高、对常见查询模式适应好。


\subsection{渐进分治生成器}

% TODO: 说明它并不直接一步到位生成 SQL，而是将问题拆为若干中间步骤
% （意图识别、过滤条件抽取、连接路径推断、聚合操作生成等）。
% 突出分步推理更适合处理复杂逻辑和长链条约束。
% 可给出一个渐进式推理过程示意图。


\section{基于执行反馈的查询修复}

% TODO: 鲁棒性增强模块。
% 列出常见错误类型：字段不存在、连接关系错误、语法错误、空结果或聚合逻辑错误。
% 说明如何把执行器反馈和原始问题一并送回模型做自反思修复。
% 交代修复轮次上限和停止条件。


\section{候选判别模型}

% 建议图表：图3-2 候选判别流程图

\subsection{训练样本构建}

% TODO: 说明判别模型的训练数据如何得到：
% 把不同生成器产生的候选 SQL 与标准 SQL 或执行结果对齐，构造成正负样本对。
% 突出"困难负例"的构造。


\subsection{配对判别与得分聚合}

% TODO: 写清为什么采用 pairwise 而不是直接多分类或绝对打分。
% 模型输入包含问题、增强 Schema 与两个候选 SQL，输出"哪个更优"的偏好判断，
% 再通过多轮比较或得分聚合选出最优候选。


\subsection{最优候选选择策略}

% TODO: 说明在线阶段的具体选择流程：锦标赛式比较、循环投票或全对比较。
% 补充是否结合执行可用性、语法合法性作为辅助信号。


\section{实验与结果分析}

\subsection{数据集构建与实验设置}

% 建议图表：表3-1 数据集统计信息表（库规模、问题数量、SQL 复杂度分布）

% TODO: 介绍自建政务数据集的"多源采集 - 逻辑注入 - 混合标注"流程。
% 给出训练集、验证集、测试集划分，以及库数量、问题数量、平均 SQL 长度、复杂查询占比等统计信息。
% 若使用 Spider、Bird 做泛化验证，也要说明使用目的和评测方式。


\subsection{评估指标与对比基线}

% TODO: 保留 EM 和 EX，以及判别准确率/F1、修复成功率、检索命中率、平均时延等指标。
% 基线分三类：传统/预训练 Text-to-SQL 模型、通用 LLM Prompt 方法、退化版本。


\subsection{总体性能对比}

% 建议图表：表3-2 不同 Text-to-SQL 方法在政务数据集上的性能对比表

% TODO: 不要只报数，要解释为什么方法在复杂政务数据集上提升更明显，
% 并区分 EM 与 EX 的含义差异。


\subsection{消融实验}

% 建议图表：表3-3 核心模块消融实验结果表

% TODO: 围绕 LA-Schema、样例库、Few-shot 生成器、分治生成器、判别模型、修复器六个组件展开。
% 从"去掉某个模块后性能下降了什么"转到"该模块解决了什么问题"。


\subsection{异构协同与判别效能分析}

% TODO: 专门证明核心创新。
% 比较"仅 Few-shot""仅分治""两者结合但无判别""两者结合加判别"的结果，
% 突出协同候选池与选择机制的贡献。


\subsection{效率分析与讨论}

% 建议图表：表3-4 效率展示分析图表

% TODO: 报告平均响应时延、候选生成数量、判别开销、修复额外耗时等指标，
% 说明性能提升是否建立在可接受的代价上。
% 至少给出单条查询平均耗时和主要耗时模块占比。


\section{本章小结}

% TODO: 归纳本章提出的核心思路：
% "用逻辑增强提高模式理解，用异构协同提高候选质量，用判别与修复提高最终鲁棒性"。
% 补充本章方法在复杂政务 SQL 生成中的效果已通过实验验证，为后续系统实现奠定基础。
